% !TEX program = xelatex
\documentclass[12pt,a4paper]{ctexart}
\usepackage{geometry}
\geometry{margin=2.4cm}
\usepackage{amsmath,amssymb}
\usepackage{booktabs}
\usepackage{enumitem}

\begin{document}

\begin{center}
{\LARGE\heiti 《药材烘干问题》审稿报告的自省总结}\\[0.5em]
{\small 对象：\texttt{report/药材烘干问题\_审稿报告.pdf}（结论：大幅修改后再审）\quad 整理日期：2026-09-12}
\end{center}

\section{一、执行原则}

对照审稿报告逐条自省后，后续修正一律遵循以下原则：

\begin{enumerate}[label=(\arabic*)]
\item \textbf{保持论文现有章节结构不变}：一至八章、5.1--5.6、6.1--6.4、7.1--7.7、8.1--8.3 的框架与小节划分不做任何调整，所有修正均落在既有段落内部（改表述、补推导、换数字），不新增或删除小节。
\item 方法/计算类问题优先修改，图表与格式类暂缓（待用户放开）。
\item 模型输入只用题给参数；外部文献数据仅用于验证对照，不进入模型。
\item 凡涉及重算的修正，改后对照题目原文、附件与 result1--4.xlsx 逐格复核，并重新编译确认。
\end{enumerate}

\section{二、对六条主要意见的自省}

\subsection{2.1 动边界模型的坐标变换与守恒关系（审稿意见 1）}

\textbf{自省：部分成立，属推导环节缺失而非结论错误。}文中假设 H8 已声明收缩各向同性且径向均匀，但未显式写出由此得到的仿射运动关系 $r_p(t)=\zeta_p R(t)$ 与 $v_r=(R'/R)r$，也未说明正是该关系使 $\zeta$ 为物质坐标、随体导数不含输运项——这是论证链条中缺失的一环。5.4 节应补 2--3 句：各向同性且径向均匀的收缩意味着任一物质点按统一比例 $R(t)/R(0)$ 收缩，故 $\zeta=r/R(t)$ 恒标记同一干物质质点，$C_t|_\zeta$ 即物质导数，输运项 $\zeta C_\zeta$ 与坐标速度项严格相消，变换后的方程不含对流项。

守恒检验部分审稿意见指出「$C$ 为干基含水率，总水量不能直接按当前体积对 $C$ 积分」。核查 7.3 节：文中已注明「守恒量按 $\sum v_iC_i$ 计，即水分浓度对体积的积分；物理水量守恒由题给有效扩散系数的标定隐含保证，见假设 H5」，与审稿意见要求的区分一致，但可以更明确：$\sum v_iC_i$ 是 $C$-方程的离散不变量（用于验证离散格式的守恒性），物理总水量应为 $\int\rho_dC\,\mathrm{d}V$（$\rho_d=\rho/(1+C)$），因题给经验式以 $C$ 为唯一状态变量，物理水量守恒由参数标定隐含保证。7.3 节补一句即可，不动结论。

\textbf{类型：}推导补充与措辞，无重算。\textbf{优先级：}第一。

\subsection{2.2 内部水分变量与空气水分变量的含义（审稿意见 2）}

\textbf{自省：成立，但整改范围应限于表述。}按题设，附件 1 给出的 $C_a$ 即「烘房水分浓度」，题面要求直接用于 Robin 边界，故模型按题给量处理 $C_a$ 为边界驱动用的等效水分浓度，无需也不应引入吸湿等温线（题给之外信息，违反只用题给参数原则）。真正的问题在 7.6 节：借 $C_a=0.04986$ 与空气饱和含湿量计算相对湿度 $\mathrm{RH}\approx57\%$，这一步隐含了「将 $C_a$ 解释为空气含湿比」的额外假设，正文未作说明。拟在 7.6 节补一句：附件 1 未指明 $C_a$ 是空气含湿比还是等效平衡水分浓度，按含湿比解释得 $\mathrm{RH}\approx57\%$；两种解释下湿球温度约 41.5 $^\circ$C、$q_{\mathrm{evap}}/q_{\mathrm{sens}}\approx3\sim5$ 的同量级结论不变。同时在符号说明中明确 $C_a$ 的题设角色。

\textbf{类型：}措辞补充。\textbf{优先级：}第一。

\subsection{2.3 「经验参数已吸收汽化潜热」的依据不足（审稿意见 3）}

\textbf{自省：成立，当前表述过强。}核查后确认需改写的断言共四处：（1）H5「其差值正体现了解吸与汽化热效应的贡献，说明经验式已隐含相变吸热的影响」；（2）2.2 节「其数值已隐含吸收了汽化与吸附等微观机制的净效应」；（3）7.6(3)「题给经验式\textbf{必须}以有效参数方式吸收潜热效应」；（4）摘要「汽化吸热的影响已被题给经验式以有效参数形式吸收，显式叠加潜热汇项将与经验参数重复计账」。$E_a=32.0$ kJ/mol 高于纯水自扩散活化能仅说明 $D$ 含解吸/汽化的综合能垒，不唯一证明热量方程已含潜热；$q_{\mathrm{evap}}/q_{\mathrm{sens}}\approx3\sim5$ 反而说明潜热可能是一阶因素。按审稿意见改写为：受题给参数限制，主模型采用显热传导与有效水分扩散系数；现有信息不足以唯一判断潜热是否已包含在经验参数中；显式叠加潜热汇项存在重复计账风险且需题给之外的参数支撑；将 $+40\%\sim50\%$ 的时长敏感性区间作为模型局限保留（正文结论处已保留，仅需把「必须吸收」改为「可能已被吸收，无法唯一判定」）。

\textbf{类型：}措辞改写（四处，均在段落内）。\textbf{优先级：}第一。

\subsection{2.4 非线性与耦合迭代缺少严格停止准则（审稿意见 4）}

\textbf{自省：成立。}当前每步固定两轮 Picard 迭代、以第三轮相对第二轮 0.2\%--0.9\% 的变化作收敛依据，不能保证各阶段达到同等精度，且 0.2\%--0.9\% 是与「单步变化」的相对量，与时间步误差量级的对比不严谨。拟在求解器中实现基于 $T$、$C$ 无穷范数相对变化的停止准则（如小于 $10^{-8}$，最大迭代 20 轮），重跑问题 2、3、4 的全部算例，以实际达到的残差更新 5.5(3) 与 7.2 节的表述；并补充说明变物性系数取迭代中最新的可用层（时间平均仍取 $n$、$n+1$ 两层算术平均，与 Crank--Nicolson 一致）。问题 1 的 $D(C)$ 仅随 $C$ 缓慢变化，且表 1、表 2 已以 Richardson 估计 $2.9\times10^{-5}$ 独立验证过精度，不改动其求解器与数据，避免无谓重算。

\textbf{类型：}计算重算，表 3--6 与 result2--4.xlsx 可能产生第四位小数的变动，须按 P4 复核。\textbf{优先级：}第一。

\subsection{2.5 烘干结束时间的数据不支持二阶 Richardson 外推（审稿意见 5）}

\textbf{自省：成立，且是本次审稿最尖锐的一条。}复核审稿意见的计算：问题 3 序列 56.7069/56.9674/57.0875 h，相邻差值比 $0.2605/0.1201\approx2.17$，实测收敛阶 $p\approx1.12$；问题 4 序列 50.5417/50.6931/50.7722 h，差值比 $0.1514/0.0791\approx1.91$，$p\approx0.94$。均非二阶，故正文「按二阶收敛 Richardson 外推 57.1275 h（±0.05 h）、50.7986 h（±0.03 h）」的误差估计依据不足，须撤换。拟采取：（1）各加一档空间离散（问题 3：$\Delta r=0.125\to0.0625$ mm，$N=321$；问题 4：$\Delta\zeta=1/80\to1/160$，$N=161$），得到四点序列后按\textbf{实测收敛阶}外推，并以最细两档的差作为保守不确定度；（2）事件时间按 $g(t)=\max_r C(r,t)-0.15$ 在相邻时间层间插值求根（现 $\Delta t=2.5$/5 s，对应误差不超过 0.001 h，影响可忽略，但按审稿意见实现）；（3）相应改写 5.6、6.3、6.4、7.2 节中「二阶外推 ±0.05 h/±0.03 h」的表述为实测收敛阶下的估计。$t^*$ 本身仍按最细离散的结果输出四位小数（题面格式要求），不确定度单独注明。

\textbf{类型：}计算重算＋表述改写。\textbf{优先级：}第一。

\subsection{2.6 问题 3 与问题 4 的差异不能全部归因于收缩（审稿意见 6）}

\textbf{自省：基本已在文中处理，仅缺对照表的一个格子。}核查 6.4 节：正文已明确「净效应由两个相反方向的因素合成」——附录 4 扩散系数较小使干燥变慢，收缩使 $D/R^2$ 增大 2.8 倍并缩短路径使干燥加快；并以附录 4 物性做了消融对照（$t_{4,F}=128.88$ h、$t_{4,S}=50.77$ h），「缩短 6.3 h（11.1\%）」的表述均以「与问题 3 对比」限定，未声称是纯粹收缩效应。审稿 2×2 表（下表）中仅缺 $t_{3,S}$（附录 3 物性＋收缩）一格：

\begin{center}
{\zihao{5}\begin{tabular}{lcc}
\toprule
物性参数 & 固定半径 & 收缩半径 \\
\midrule
附录 3 & $t_{3,F}=57.09$ h（已有） & $t_{3,S}=?$（待补） \\
附录 4 & $t_{4,F}=128.88$ h（已有） & $t_{4,S}=50.77$ h（已有） \\
\bottomrule
\end{tabular}}
\end{center}

补跑附录 3 物性＋附件 2 收缩的算例即可闭合，使「收缩在两种物性下均加速干燥」或「收缩效应依赖物性」有直接证据。

\textbf{类型：}补一个算例（成本低）。\textbf{优先级：}第一。

\section{三、对十二条其他意见的自省}

\begin{center}
{\zihao{5}\begin{tabular}{p{0.55\textwidth}p{0.35\textwidth}}
\toprule
意见 & 自省结论与响应 \\
\midrule
1. 轴向/径向扩散时间比应为 $(L/(2R))^2\approx39$ 而非 156 & 成立。2.1 节端面特征尺度取半长 $L/2$，比值 156 改为 39，结论（轴向可忽略）不变。措辞修正，第一优先级 \\
2. $D/h_m$ 宜称 Robin 边界特征长度（外推长度），扩散深度用 $\sqrt{Dt}$ & 成立。6.1、5.6 的「边界层厚度」改「特征长度（外推长度）」，$\sqrt{Dt}$ 估算 6.1 节已采用。措辞修正 \\
3. 变物性热方程应说明为弱耦合近似 & 基本已符合：2.2 节已称「双向弱耦合」，5.2 节「双向耦合」补「弱」并加一句忽略水分输运携带显焓的说明。措辞修正 \\
4. 2--3 天只能作合理性检查 & 已符合：正文用「一致」「佐证」，未作独立验证声称。无需改动 \\
5. 14400 s 后固定附件末值是重要外推假设 & 已基本符合：H7 已列为假设，7.7 节灵敏度已含 $T_a\pm2\ ^\circ$C、$C_a\pm20\%$。可在 H7 补「该外推是对题给信息的重要延伸」一句。可选 \\
6. 四位小数要求误差低于 $0.5\times10^{-4}$ & 成立。Q2（0.25 mm，估计 $6.1\times10^{-5}$）与 Q4（1/80，$5.4\times10^{-5}$）略高于阈值，按 2.5 节的加密计划一并解决（Q2 加至 0.125 mm、Q4 加至 1/160），达标后再宣称四位小数可靠。计算重算，第一优先级 \\
7. 图 1 信息过密，建议拆分 & 图表类，按既定约束暂缓 \\
8. 图 2 图题不应仅写「问题 4 的收缩模型」 & 图表类，暂缓；现图题已改述机制，若图内确为 Q3/Q4 对比，届时一并改 \\
9. 图中 56.97 h、50.69 h 与正文最终值不一致 & 图表类，暂缓；smart-charts 技能清单中已记为已知待办 \\
10. 摘要中数值方法与潜热辩护篇幅过长 & 部分成立。拟与意见 3 的潜热改写同步精简该段，保留模型、四问结果、关键验证与主要局限，维持摘要仍占满且不超出第 1 页。文字类，第二优先级 \\
11. AI 工具使用说明应写明使用环节与人工核验 & 该节为用户自填内容，需用户本人补充，我不代写 \\
12. 逐条核对参考文献 & 正文 19 条文献与各论点的对应关系上轮已审计；书目元数据（作者/题名/卷期页码）建议逐条对原刊再核一遍。低优先级待办 \\
\bottomrule
\end{tabular}}
\end{center}

\section{四、计算任务与执行顺序}

\textbf{第一轮（纯文字修正，不动数字）}：2.1 节 156→39；6.1/5.6 节「边界层」措辞；5.2 节弱耦合限定；5.4(1) 补仿射运动推导；H5、2.2、7.6(1)(3)、摘要的潜热表述软化；7.3 节补守恒量说明；7.6(1) 补 $C_a$ 解释限定。改后重编译。

\textbf{第二轮（计算重算）}：实现 Picard 停止准则并重跑问题 2--4；Q2 加密至 0.125 mm、Q4 加密至 1/160；Q3、Q4 各加一档空间离散重估 $t^*$ 收敛阶；补跑 $t_{3,S}$ 算例闭合 2×2 表。随后更新表 3--6、result2--4.xlsx 及正文中受影响的数字（含六章开头 $t^*$ 与 5.5(3)、5.6、7.2、6.4 各处的收敛数据），逐格对照 xlsx 复核后重编译。

\textbf{第三轮（待用户放开）}：图 1--3 与 smart-charts 同步更新；摘要最终精简；文献元数据核对。

\textbf{保持不变的事项}：章节结构与小节划分；问题 1 的求解器、表 1/2 与 result1.xlsx（精度已独立达标）；模型输入参数（仅题给）；题目与附件数据。

\section{五、自省结论}

审稿报告的六条主要意见中，第 3、4、5 条指出了当前版本的真实缺陷（潜热表述过强、迭代无停止准则、$t^*$ 二阶外推不成立），必须修改；第 1、2 条是论证链缺失与表述问题，补推导与说明即可，模型与结论不动；第 6 条正文已基本处理，补一个算例即闭合。十二条其他意见中 8 条已处理或仅需措辞修正，4 条属图表与用户自填项暂缓。整体判断与审稿结论一致的方向是：核心物理闭环（动边界、变量定义、潜热）与事件时间收敛需要在当前结构内加强，但不需要推倒重来。

\end{document}
